% Otomatik Uretilen IEEE SIU / Bitirme Tezi Tablosu (LaTeX)
% Yazar: Ahmet Eren Cokcoskun (TA7CAE)
\begin{table}[htbp]
\centering
\caption{Kayseri Orta Enlem Iyonosferinde Yaklasim vs. Tam Magneto-Iyonik Zayiflama Analizi}
\begin{tabular}{ccccccc}
\hline
\textbf{Frekans (MHz)} & \textbf{Katman} & \textbf{$h_{\mathrm{refl}}$ (km)} & \textbf{Basit $1/f^2$ (dB)} & \textbf{O-Modu (dB)} & \textbf{X-Modu (dB)} & \textbf{Hata Payi (dB)} \\
\hline
1.840 & E / Es & 95.0 & 265.31 & 157.51 & 116.53 & 148.78 \\
3.573 & E / Es & 96.5 & 99.23 & 64.98 & 168.46 & 69.23 \\
7.074 & E / Es & 100.0 & 34.06 & 26.09 & 43.71 & 9.65 \\
10.136 & E / Es & 103.0 & 19.71 & 16.43 & 23.70 & 3.99 \\
14.074 & F1 & 168.5 & 13.91 & 12.29 & 16.68 & 2.77 \\
21.074 & F2 & 233.0 & 6.29 & 5.38 & 6.52 & 0.23 \\
\hline
\end{tabular}
\end{table}

% Gercek Irtifa Iyonogram Tersinimi (True-Height Inversion foF2, hmF2) Bulgulari
\begin{table}[htbp]
\centering
\caption{Dikey FMCW Iyonosond Iyonogramindan Sayisal Tersinim ile Kestirilen Iyonosferik Tepe Parametreleri}
\begin{tabular}{ccccc}
\hline
\textbf{$f_oF2$ (MHz)} & \textbf{$h_mF2$ (km)} & \textbf{$N_mF2$ ($\mathrm{m}^{-3}$)} & \textbf{$\mathrm{TEC}_{\mathrm{bottom}}$ (TECU)} & \textbf{Kalinti RMS (km)} \\
\hline
10.871 & 262.4 & 1.466e+12 & 10.42 & 0.0000 \\
\hline
\end{tabular}
\end{table}
